\section{Дифференциальные исчисления в линейных нормированных пространствах}

\begin{note}
	Если $A \colon E_1 \to E_2$ --- линейные оператор между линейно нормированными пространствами, то принято опускать скобки у аргумента:
	\[
		\forall x \in E_1\ \ Ax := A(x)
	\]
\end{note}

\begin{note}
	Далее, если не оговорено обратного, мы считаем $E_1$ и $E_2$ всегда линейно нормированными пространствами. Также, это может явно не проговариваться, но иногда мы исключаем случай $E_1$ из доказательства, благо он достаточно тривиален, чтобы это было уместно.
\end{note}

\begin{definition}
    Линейный оператор $A \colon E_1 \to E_2$ называется \textit{ограниченным}, если 
    \[
        \exists M \in \R \such \forall x \in E_1\ \ \|Ax\|_{E_2} \leq M
    \]
\end{definition}

\begin{theorem}
	Пусть $A \colon E_1 \to E_2$ --- линейный оператор между ЛНП. Тогда он ограничен $\lra$ он непрерывен (в любой точке $E_1$).
\end{theorem}

\begin{proof}~
    \begin{itemize}
        \item Пусть $A$ --- ограниченный линейный оператор: $E_1 \to E_2$. Из ограниченности есть такая оценка:
        \[
            \exists M > 0 \such \|Ax_1 - Ax_2\|_{E_2} = \|A(x_1 - x_2)\|_{E_2} \leq M\|x_1 - x_2\|_{E_1}
        \]
        Отсюда сразу следует непрерывность:
        \[
            \forall \eps > 0\ \exists \delta := \frac{\eps}{M} > 0 \such \forall x_1, x_2 \in E_1,\  \|x_1 - x_2\|_{E_1} < \delta \emb \|Ax_1 - Ax_2\|_{E_2} < \eps
        \]
        
        \item Непрерывность во всех точках означает как минимум непрерывность в нуле:
        \[
            \forall \eps > 0\ \exists \delta > 0 \such \forall h \in E_1, \|h\|_{E_1} < \delta \quad \|Ah\|_{E_2} < \eps
        \] 
        Зафиксируем $\eps = 1$ и, следовательно, $\delta > 0$. Тогда, для любого ненулевого $x \in E_1$ можно выбрать $h = (\delta / 2) \cdot x / \|x\|_{E_1}$, удовлетворяющий условию непрерывности в нуле. Подставим $h$:
        \[
        	\|Ah\|_{E_2} = \nr{A \frac{\delta x}{2\|x\|_{E_1}}}_{E_2} = \nr{\frac{\delta}{2\|x\|_{E_1}}Ax}_{E_2} < \eps
        \]
        Так как $A$ --- линейный оператор, то положительную константу можно вытащить за знак нормы:
        \[
              \nr{\frac{\delta}{2\|x\|_{E_1}}Ax}_{E_2} = \frac{\delta}{2\|x\|_{E_1}}\|Ax\|_{E_2} < \eps = 1 \Lora \|Ax\|_{E_2} < \frac{2}{\delta}\|x\|_{E_1}
        \]
        Следовательно, искомая константа $M = 2 / \delta$.
    \end{itemize}
\end{proof}

\begin{definition}
    Пусть $A \colon E_1 \to E_2$ --- линейный ограниченный оператор между ЛНП. \textit{Нормой оператора} $A$, обозначаемой как $\|A\|$, называется следующая величина:
    \[
    	\|A\| := \inf \{M \in \R \colon \forall x \in E_1\ \|Ax\|_{E_2} \le M\|x\|_{E_1}\}
    \]
\end{definition}

\begin{proposition}
    Если $A \colon E_1 \to E_2$ --- линейный ограниченный оператор, то имеет место формула для нормы (за исключением тривиального случая $E_1 = 0$, в котором все операторы ограничены и имеют норму $0$):
    \[
        \|A\| = \sup_{x \in E_1,\ \|x\| = 1} \|Ax\|_{E_2} = \sup_{x \in E,\ x \neq 0} \frac{\|Ax\|_{E_2}}{\|x\|_{E_1}} 
    \]
\end{proposition}

\begin{proof}
	\begin{enumerate}
		\item Установим равенство нормы с последним супремумом. Для этого поймём, почему оно вообще может быть верным. Если в определении нормы $x \neq 0$, то
		\[
			\|Ax\|_{E_2} \le M\|x\|_{E_1} \lra \frac{\|Ax\|_{E_2}}{\|x\|_{E_1}} \le M
		\]
		Стало быть, имеет место общее неравенство
		\[
			\sup_{x \in E_1,\ x \neq 0} \frac{\|Ax\|_{E_2}}{\|x\|_{E_1}} \le \|A\|
		\]
		Предположим противное: неравенство выполнилось строго. Тогда $\exists M$ такое, что оно лежит в интервале между супремумом и нормой:
		\[
			\exists M \in \R \such M \in \ps{\sup_{x \in E_1,\ x \neq 0} \frac{\|Ax\|_{E_2}}{\|x\|_{E_1}}; \|A\|}
		\]
		Это сразу же приводит к противоречию, ибо из факта $M < \|A\|$ следует, что для всех $x \neq 0$ выполнено неравенство из определения, но и для нуля оно верно автоматически (ибо $A(0) = 0$).
		
		\item Установим равенство супремумов. Действительно:
		\[
			\forall x \in E,\ x \neq 0\ \exists h := \frac{x}{\|x\|},\ \|h\| = 1
		\]
		Более того, имеет место такое соотношение:
		\[
			\frac{\|Ax\|_{E_2}}{\|x\|_{E_1}} = \frac{\|x\|_{E_1} \cdot \|Ah\|_{E_2}}{\|x\|_{E_1}}
		\]
		Иначе говоря, каждому элементу множества $\|Ax\|_{E_2} / \|x\|_{E_1}$ соответствует какой-то элемент $\|Ah\|_{E_2}$ и наоборот. Значит, супремумы этих множеств будут совпадать.
	\end{enumerate}
\end{proof}

\begin{proposition}
	Если $\Lin(E_1, E_2)$ --- множество линейных ограниченных операторов из $E_1$ в $E_2$, то оно является линейно нормированным пространством с нормой $\|A\|$.
\end{proposition}

\begin{proof}
	Проверим 3 аксиомы нормы:
	\begin{enumerate}
		\item Если $\|A\| = 0$, то $\forall x \in E_1\ \|Ax\|_{E_2} = 0$. Следовательно, $A = 0$ --- нулевой оператор.
		
		\item Определим операцию умножения оператора $A$ на константу $\alpha \in \R$ поточечно:
		\[
			\forall x \in E_1 \quad (\alpha A)(x) := \alpha \cdot A(x)
		\]
		В силу доказанной формулы нормы оператора, вынесение константы за знак нормы оператора работает как надо.
		
		\item Проверим, что выполняется неравенство треугольника:
		\begin{multline*}
			\|A_1 + A_2\| = \sup_{x \in E_1,\ \|x\|_{E_1} = 1} \|(A_1 + A_2)x\|_{E_2} \leq
			\\
			\sup_{x \in E_1,\ \|x\|_{E_1} = 1} \|A_1x\|_{E_2} + \sup_{x \in E_1,\ \|x\|_{E_1} = 1} \|A_2x\|_{E_2} = \|A_1\| + \|A_2\|
		\end{multline*}
	\end{enumerate}
\end{proof}